% 仙后座（综述）
% license CCBYSA3
% type Wiki

本文根据 CC-BY-SA 协议转载翻译自维基百科\href{https://en.wikipedia.org/wiki/Cassiopeia_(constellation)}{相关文章}。

仙后座（Cassiopeia，ⓘ）是位于北天的一個星座與星群，其名稱來自希臘神話中虛榮的王后卡西奧佩婭（Cassiopeia），她是安德洛墨達（Andromeda）的母親，曾誇耀自己擁有無與倫比的美貌。仙后座是公元 2 世紀希臘天文學家托勒密列出的 48 個星座之一，至今仍是現代 88 個星座之一。由於其由五顆亮星組成的獨特「W」形狀，仙后座極易辨認。

仙后座位於北天，在北緯 34° 以上的地區全年可見；在（亞）熱帶地區，9 月至 11 月初觀測條件最佳；而在南半球低緯度、低於南緯 25° 的熱帶地區，則可於特定季節在北方低空看到。

視星等為 2.2 的仙后座 α 星，即舍達爾（Schedar），是仙后座中最亮的恆星。該星座擁有一些已知最為明亮的恆星，包括黃色特超巨星仙后座 ρ 星（Rho Cassiopeiae）與 V509 Cassiopeiae，以及白色特超巨星仙后座 6 星。1572 年，第谷·布拉赫（Tycho Brahe）所觀測到的超新星就在仙后座中明亮爆發$^{[5]}$。仙后座 A 是一個超新星殘骸，在頻率高於 1 GHz 時，是全天最亮的銀河系外射電源。已有 14 個恆星系統被發現擁有系外行星，其中之一——HD 219134——被認為可能擁有六顆行星。一段富集的銀河帶穿過仙后座，其中包含多個疏散星團、年輕而明亮的銀河盤恆星以及星雲。IC 10 是一個不規則星系，是目前已知距離最近的星暴星系，也是本星系群中唯一的此類星系。
\subsection{神話}
\begin{figure}[ht]
\centering
\includegraphics[width=6cm]{./figures/99da126f8a82dd17.png}
\caption{《乌拉尼娅之镜》中描绘的坐在宝座上的仙后座（Cassiopeia）} \label{fig_Cape_1}
\end{figure}
该星座以埃塞俄比亚王后卡西奥佩娅（Cassiopeia）命名。卡西奥佩娅是埃塞俄比亚国王刻甫斯（Cepheus）的妻子$^{[5]}$，也是公主安德洛墨达（Andromeda）的母亲。刻甫斯与卡西奥佩娅以及安德洛墨达一同被安置在星空中，相互毗邻。她之所以被放置在天空中，作为惩罚，是因为她夸口说自己的女儿安德洛墨达比涅瑞伊得斯（Nereids，海中女神）更加美丽，或者另一种说法是，她自称比海中仙女本身更加美丽，从而激怒了波塞冬$^{[6]}$。她被迫坐在宝座上绕北天极旋转，一半时间必须紧紧抓住宝座以免跌落；而波塞冬则裁定安德洛墨达应被绑在岩石上，作为海怪刻托斯（Cetus）的祭品。随后，安德洛墨达被英雄珀尔修斯所拯救，并最终与其成婚$^{[7][8]}$。

在历史上，卡西奥佩娅星座的形象有多种不同的描绘方式。在波斯，她被苏菲（al-Sufi）描绘为一位女王，右手持有一根带新月的权杖，头戴王冠，身旁还有一只双峰骆驼；在法国，她则被描绘为坐在大理石宝座上，左手持棕榈叶，右手提着自己的长袍。这一形象出自奥古斯丁·罗耶（Augustin Royer）于1679年出版的星图集$^{[7]}$。

在中国天文学中，构成仙后座的恒星分布于三个区域之中：紫微垣（紫微垣，Zǐ Wēi Yuán）、北方玄武（北方玄武，Běi Fāng Xuán Wǔ）以及西方白虎（西方白虎，Xī Fāng Bái Hǔ）。

中国天文学家在现代所称的仙后座区域内辨认出多种星象组合。κ、η 与 μ 仙后座星组成了名为“王桥”的星官；当它们与 α 与 β 仙后座星共同观测时，则构成了名为“王良”的大车。御者手中的鞭子由 γ 仙后座星表示，该星有时被称为“次”，即汉语中“鞭”的意思$^{[7]}$。

在印度神话中，仙后座（Cassiopeia）与神话人物沙尔米什塔（Sharmishtha）相关联——她是伟大的阿修罗（Daitya）国王弗里什帕尔瓦（Vrishparva）的女儿，也是德瓦雅尼（Devayani，对应安德洛墨达）的朋友。

在威尔士神话中，Llys Dôn（字面意思为“多恩的宫廷”）是该星座的传统威尔士名称。多恩的至少三位子女同样具有天文学上的对应：Caer Gwydion（“圭迪恩的堡垒”）是银河的传统威尔士名称，而 Caer Arianrhod（“阿丽安罗德的堡垒”）则对应北冕座（Corona Borealis）$^{[9]}$。

在17世纪，仙后座中的恒星曾被用来描绘多位《圣经》人物，其中包括所罗门之母拔示巴（Bathsheba）、《旧约》中的女先知底波拉（Deborah），以及耶稣的追随者抹大拉的玛利亚（Mary Magdalene）$^{[7]}$。

在一些阿拉伯星图集中，仙后座的恒星还被描绘为一个名为“染色之手”（Tinted Hand）的形象。据不同说法，这只手要么象征一只被指甲花染成红色的女子之手，要么象征穆罕默德之女法蒂玛（Fatima）沾满鲜血的手。该“手”由 α、β、γ、δ、ε 与 η 仙后座星组成；其“手臂”则由 α、γ、δ、ε、η 与 ν 英仙座星组成$^{[7]}$。

另一个包含仙后座恒星的阿拉伯星座是“骆驼”。其头部由仙女座的 λ、κ、ι 与 φ 星构成；驼峰是 β 仙后座星；身体由其余仙后座恒星组成；而双腿则由英仙座与仙女座中的恒星构成$^{[7]}$。

在其他文化中，人们还从这一星群中看到了手或麋鹿鹿角的形象$^{[10]}$。例如在萨米人（Sámi）的传统中，仙后座的 “W” 形状被视为麋鹿的鹿角；西伯利亚的楚科奇人（Chukchi）同样将其中五颗主星视作五只驯鹿雄鹿$^{[7]}$。

马绍尔群岛的居民将仙后座视为一幅巨大鼠海豚星座的一部分：仙后座的主星构成其尾部，仙女座与三角座（Triangulum）构成其身体，而白羊座则构成其头部$^{[7]}$。在夏威夷，α、β 与 γ 仙后座星各自有名称：α 仙后座星被称为 Poloahilani，β 仙后座星被称为 Polula，γ 仙后座星被称为 Mulehu。普卡普卡（Pukapuka）群岛的人们则将仙后座视为一个独立的星座，称为 Na Taki-tolu-a-Mataliki $^{[11]}$。
\subsection{特征}
\begin{figure}[ht]
\centering
\includegraphics[width=6cm]{./figures/8d913297d824a2aa.png}
\caption{夜空中的仙后座} \label{fig_Cape_2}
\end{figure}